% ARCHIVED DRAFT: independently audited by root and incorporated on 2026-11-19 UTC.
% The assembled paper is authoritative; this file retains drafting provenance.
\begin{saresult}{B.073}{Proved under stated hypotheses}{Zero-dimensional bases with T-nilpotent nilradical}
Let $R$ be a nonzero commutative ring with nilradical $N$. Suppose
$R/N$ is absolutely flat and $N$ is T-nilpotent: for every sequence
$n_1,n_2,\ldots$ in $N$, some finite product $n_1\cdots n_j$ is zero.
Then every injective ring epimorphism out of $R$ is an isomorphism.
Consequently, for every unital inclusion $R\subset S$, the following
are equivalent:
\[
\begin{gathered}
 \text{the polynomial co-operations exist uniquely};\\
 \theta_2\text{ is injective};\qquad C=R;\\
 (R:_S D_d(R))=R\quad\text{for every }d\ge1;\\
 \operatorname{Ann}_{S/R}(D_d(R))=0\quad\text{for every }d\ge1;\\
 \Omega(R,S)=R[X].                                       \tag{B73.1}
\end{gathered}
\]
All positive cases have every multivariable polynomial tensor
isomorphism. In particular the theorem applies to every zero-dimensional
ring with T-nilpotent nilradical; the nilradical need not be nilpotent.
\end{saresult}
\begin{proof}
Let $R\subset B$ be a ring epimorphism. The ideal $NB$ is nil: each of
its elements is a finite sum of multiples of nilpotents in a commutative
ring, and is therefore nilpotent. Thus $NB\cap R=N$, and base change
gives an injective epimorphism
\[
                         R/N\longrightarrow B/NB.
\]
The absolutely flat argument in B.065 makes this an isomorphism.
Consequently $B=R+NB$.

For completeness we prove that T-nilpotence suffices to conclude
$B=R$, without imposing finite generation on $B$ or $N$. Fix $b\in B$.
Write $b=r+\sum_{i=1}^s n_i b_i$ with $r\in R$, $n_i\in N$, and repeat
this finite expansion for every $b_i$, then for each new coefficient,
and so forth. This gives a rooted, finitely branching tree. At a node
the coefficient of its remaining $B$-element is the product of the
successive $N$-labels along the path from the root. Remove nodes where
that product is zero. There is no infinite path in the remaining tree,
by T-nilpotence. Such a finitely branching tree is finite: if it had
infinitely many nodes, the root would have a child with infinitely many
descendants, and repeating that finite-choice argument would give an
infinite path. Hence after a finite expansion all remaining coefficients
are zero and the displayed expression for $b$ lies in $R$. Thus $B=R$.

If $\theta_2$ is injective, B.064 applies this fact to the actual
coefficient ring $C$ and gives $C=R$. The converse and every equivalent
ideal condition in (B73.1) follow from B.064. The positive tensor maps
are ordinary polynomial tensor maps.

Finally a reduced zero-dimensional commutative ring is absolutely flat;
this is the equivalence in the opened Stacks Lemma 15.106.5, tag 092F,
already used in B.065. Thus the stated zero-dimensional class satisfies
the hypotheses. The T-nilpotence requirement is substantive: in B.067,
a sequence of degree-zero monomials from distinct tree components has
every finite product nonzero, so that example lies outside this class.
\end{proof}
